% chapters/abstract-en.tex - English Abstract

\chapter*{Abstract}
\addcontentsline{toc}{chapter}{Abstract}

\textbf{Objective:} Given that standardized repetitions-in-reserve (RIR)-based self-guidance has established scientific foundation and accessibility, empirical evidence remains limited regarding whether a computer-vision-based mobile digital monitoring solution provides incremental adaptive benefit and implementation feasibility justifying its additional cost. This study employed a three-phase sequential design to systematically evaluate the measurement foundation, intervention feasibility, exploratory training effects, and user acceptance of a self-developed mobile computer vision system (SportSci Pro) for digital monitoring of lower-body resistance training in young men.

\textbf{Methods:} Study 1 recruited 13 healthy young men to synchronously collect mean concentric velocity (MCV) data during barbell back squats using SportSci Pro and GymAware under controlled incremental-load conditions (60 paired observations). Study 2 employed a two-arm parallel, open-label, stratified block-randomized controlled pilot design with 36 randomized participants (AI-assisted $n$ = 18, self-guidance $n$ = 18) undergoing a 4-week, twice-weekly lower-body resistance training intervention.

\textbf{Results:} Study 1 yielded ICC = 0.940 (MAE = 0.047 m/s); real-world rep-level paired analysis yielded ICC = 0.991. In the PP sample, training completion rate reached 98.4\%, AI group velocity target hit rate reached 80.2\%, and between-group difference in total repetitions was $-$5.19 ($p$ = 0.035). For CMJ, adjusted difference was +2.16 cm ($g$ = 0.85); for training self-efficacy, difference was +9.88 points ($p$ $<$ 0.001).

\textbf{Conclusion:} SportSci Pro provides a measurement foundation for deploying low-cost objective velocity feedback in general population resistance training. The 4-week intervention did not produce confirmatory between-group superiority, but captured a directionally consistent signal in CMJ and produced a significant advantage in training self-efficacy. The current protocol should be positioned as a training support tool with application potential that requires large-sample, long-term validation.

\addcontentsline{toc}{chapter}{Abstract}

\textbf{Key Words:} digital monitoring; resistance training; velocity-based training; pilot study; mixed methods research; randomized controlled pilot trial
